\section{Compilation Strategy: Deterministic Unrolling}\label{sec:compilation-strategy}

Motivo is designed as a fully compiled language, distinguishing it from interpreted musical environments.
This design choice is driven by the mandate for a \textit{zero-runtime architecture}, where all musical and logical
structures are resolved before any sound is produced or any output file is written.

\subsection{The Zero-Runtime Mandate}\label{detail-subsec:the-zero-runtime-mandate}

In a zero-runtime system, the complexity of the composition is handled entirely by the compiler.
This means that features such as nested loops, conditional branching, and recursive pattern calls do not exist in the
output artifact.
Instead, the compiler performs a process of \textit{deterministic unrolling}.
Every loop is expanded into its constituent iterations, every conditional branch is evaluated and resolved, and every
pattern call is inlined.

The result of this process is an Intermediate Representation (IR) that consists of a flat, linear list of atomic musical
events.
This architecture provides several benefits:

\noindent \textbf{High Performance.} The output artifact is extremely lightweight, as all logical overhead has been
removed.

\noindent \textbf{Offline Validation.} Semantic errors---such as variable shadowing, type mismatches, or
invalid rhythmic
math---are caught during the compilation phase, providing immediate feedback to the composer.

\noindent \textbf{Portability.} The resulting IR can be easily serialized into formats beyond MIDI in future extensions
without requiring a specialized interpreter to read it.

\subsection{Temporal Reconstruction of Patterns}\label{detail-subsec:temporal-reconstruction-of-patterns}

A unique aspect of the compilation strategy is how the temporal span of a pattern is reconstructed.
When a pattern is called, the compiler does not just emit the notes within it; it calculates the total duration occupied
by the pattern, including internal gaps and trailing rests.
This is achieved by capturing the state of the internal cursor before and after the pattern's body is unrolled.
The difference between these two points defines the pattern's ``footprint'' on the timeline, allowing the pattern to be
treated as a single, opaque musical unit by the caller.

\subsection{Static Verification of Compositional Logic}\label{detail-subsec:static-verification-of-compositional-logic}

Because the language is compiled, it can enforce constraints that would be difficult to manage in an interpreted system.
The semantic analysis pass verifies that the composition adheres to the physical limits of the target protocol (e.g.,
ensuring notes are within the MIDI pitch range of 0--127) and that the structural hierarchy is valid (e.g., preventing a
\texttt{track} from being defined inside a \texttt{pattern}).
This validation is intended to ensure that the final MIDI artifact is not only syntactically correct but also musically
meaningful and interoperable with MIDI-aware tools.

\begin{figure}[htbp]
  \centering
  \begin{tikzpicture}[
      box/.style={rounded corners=4pt, draw=motivoInk!60, very thick, minimum width=2.75cm, minimum height=0.95cm,
      align=center, font=\footnotesize\sffamily},
      arrow/.style={-Stealth, very thick, draw=motivoInk!75}
    ]
    \node[box, fill=motivoBlue!14] (ast) {AST statement\\\texttt{loop (500)}};
    \node[box, fill=motivoCyan!18, right=0.55cm of ast] (exec) {Lowering pass\\evaluate loop body};
    \node[box, fill=motivoGreen!18, right=0.55cm of exec] (ir) {IR events\\absolute beats};
    \node[box, fill=motivoAmber!20, right=0.55cm of ir] (smf) {SMF events\\delta-times};

    \draw[arrow] (ast) -- node[above, font=\scriptsize\sffamily] {control flow} (exec);
    \draw[arrow] (exec) -- node[above, font=\scriptsize\sffamily] {note events} (ir);
    \draw[arrow] (ir) -- node[above, font=\scriptsize\sffamily] {serialization} (smf);

    \node[below=0.55cm of exec, align=center, font=\small\sffamily, text=motivoInk] {
      The output no longer contains \texttt{loop}, \texttt{for}, or \texttt{if}.\\
      It contains timestamped note events that the MIDI writer can sort and encode.
    };
  \end{tikzpicture}
  \caption{Lowering removes source-level control flow by expanding it into concrete musical
  events.}\label{fig:ast-to-ir}
\end{figure}
